\setchapterpreamble[u]{\margintoc}
\chapter{Достижимость с ограничениями в виде контекстно-свободных языков}
\label{chpt:CFPQ_Overview}
\tikzsetfigurename{CFPQ_Overview_}

В данной главе мы дадим обзор задачи достижимости с ограничениями в виде контекстно-свободных языков.
Будут рассмотрены история вопроса, области применения, уточнена формальная постановка для КС-случая и описана общая идея адаптации классических алгоритмов синтаксического анализа к графам.

\section{История вопроса}

TODO: История задачи КС-достижимости. Первые работы Яннакакиса~\sidecite{Yannakakis}, Репса~\sidecite{Reps} и других.

\section{Области применения}

TODO: Статический анализ кода (межпроцедурный анализ, анализ указателей, анализ алиасов, срезы программ), графовые базы данных, биоинформатика.
Часть материала может быть перенесена из главы~\ref{chpt:FLPQ} (§«Области применения»).

\section{Уточнение формальной постановки для КС-случая}

TODO: Пусть дан граф $\mscrG = \langle V, E, L \rangle$ и КС-грамматика $G = \langle \Sigma, N, P, S \rangle$, где $L \subseteq \Sigma$.
Задача достижимости: найти множество пар вершин $(v_i, v_j)$, для которых существует путь $\pi$, такой что $\omega(\pi) \in L(G)$.
Альтернативно: для заданных множеств стартовых и финальных вершин.

\section{Общая идея адаптации алгоритмов синтаксического анализа к графам}

TODO: Ключевая идея~--- замена позиций в строке на вершины графа.
При обработке терминала в строке есть ровно один следующий символ; в графе~--- несколько исходящих рёбер (конфликт типа shift-shift).
Алгоритмы должны обрабатывать все возможные продолжения.
Различные стратегии обработки этого конфликта приводят к различным алгоритмам КС-достижимости (гл.~\ref{chpt:CFPQ_CYK}--\ref{chpt:GLR}, а также гл.~\ref{chpt:MatrixBasedAlgos}--\ref{chpt:MCFLPQ}).

\section{Структура дальнейшего изложения}

TODO: Описать, в каком порядке будут рассмотрены алгоритмы КС-достижимости в последующих главах.
